% ==================== SECTION CONVERTISSEUR D'UNITES ====================

\element{convertisseur}{
    \section{Convertisseur d'unités électriques}
    Un technicien souhaite créer un outil de conversion pour les unités électriques couramment utilisées en électronique.

    \begin{question}{convert-1}
        Ecrire une fonction \lstinline[language=C]|float milli_vers_unite(float valeur)| qui convertit une valeur en milliunites (mA, mV, mW...) vers l'unite de base (A, V, W...).

        Par exemple : \lstinline[language=C]|milli_vers_unite(250.0)| doit retourner \lstinline[language=C]|0.25|

        Ecrire egalement un petit programme \lstinline[language=C]|main| qui teste cette fonction en demandant une valeur a l'utilisateur et en affichant le resultat de la conversion.

        \textbf{Pour cette question, écrire un programme complet (avec les \#include nécessaires).}

        \evaluationProfGrille{2}{15}
    \end{question}

    \begin{question}{convert-2}
        Écrire une fonction plus complète qui calcule la puissance électrique et la convertit dans l'unité demandée :

        \lstinline[language=C]|float convertir_puissance(float tension, float courant, char unite)|

        Cette fonction doit :
        \begin{itemize}
            \item Calculer la puissance : $P = U \times I$
            \item Le paramètre \lstinline[language=C]|unite| indique l'unité de sortie souhaitée :
        \begin{itemize}
            \item \texttt{W} : retourner la puissance en Watts
            \item \texttt{m} : retourner la puissance en milliwatts (mW)
            \item \texttt{k} : retourner la puissance en kilowatts (kW)
        \end{itemize}
            \item Pour tout autre caractère : retourner la puissance en Watts
        \end{itemize}

        \textbf{Écrire uniquement la fonction (pas besoin du main ni des \#include).}

        \evaluationProfGrille{2}{12}
    \end{question}

    \begin{question}{convert-3}
        Écrire une fonction \lstinline[language=C]|void afficher_table_conversion(float puissance_w)| qui affiche la table de conversion complète d'une valeur de puissance donnée en Watts.

        La fonction doit afficher les conversions principales depuis les Mégawatts (MW) jusqu'aux milliwatts (mW). Un exemple d'affichage pour une puissance de 1500 W est donné ci-dessous :
        \lstinputlisting[style=console]{sources/codes/exemple-conversion.txt}


        \textbf{Consignes importantes :}
        \begin{itemize}
            \item Utiliser une boucle pour générer les conversions
            \item Formater l'affichage avec 6 décimales
            \item La puissance en Watts est toujours la référence et la puissance de 10 doit être indiquée comme dans l'exemple
        \end{itemize}

        \textbf{Écrire uniquement la fonction (pas besoin du main ni des \#include).}

        \evaluationProfGrille{2}{12}
    \end{question}
}
